% Resume section entries for OpenAI application
% Copy into your LaTeX resume (e.g. under Experience or Projects)

% --- Sniff ---
\textbf{Developer, Sniff (TreeHacks 2026)} \hfill Feb 2026\\
\textit{Hackathon project; LLM-powered shopping safety agent}
\begin{itemize}
  \item Built a multi-turn LLM agent that clarifies user intent, searches multi-retailer listings, and runs 5 parallel fraud checks (WHOIS, Safe Browsing, ScamAdviser, Reddit, brand impersonation) with weighted trust scoring and live streaming via SSE.
  \item Implemented tool-orchestration pipeline with function calling and real-time result streaming; frontend in Next.js/React with structured fraud reports and price-ranked recommendations.
\end{itemize}

% --- AI-Kotoba ---
\textbf{Developer, AI-Kotoba} \hfill 2025 -- Present\\
\textit{Personal project; macOS app for learning Japanese via AI conversations}
\begin{itemize}
  \item Built a native macOS app (SwiftUI, SwiftData) that generates scenario-based Japanese conversations with Chinese translations using a two-turn prompting strategy (monolingual generation then translation) via Claude/OpenAI APIs.
  \item Implemented speech-to-text role-play practice with AI feedback, TTS pronunciation, vocabulary management, and SM-2 spaced-repetition flashcards; API keys stored in Keychain.
\end{itemize}
